\section{Methodology}
\label{sec:method}

\subsection{Overview}

The study is a structured content analysis of privacy communication, combined
with a task-based inspection of the interfaces that carry part of that
communication. It proceeds in four steps: build a frozen sample of apps; collect
every privacy-relevant text those apps put in front of a German user; walk
through a stratified subsample of the apps manually to capture what the
interface says in use; then code the resulting corpus against a harm-oriented
framework, with LLM assistance for volume and manual double-coding for
validation. Figure~\ref{fig:method} shows the pipeline.

Two design commitments run through all four steps. First, \textbf{every code
must be anchored to a verbatim quote} with its source, location, and capture
date; a code without a quote does not enter the results. Second, \textbf{every
instrument is versioned} --- the sampling protocol, the codebook, and each LLM
prompt --- so that the analysis is reproducible from the repository alone.

\begin{figure}[htbp]
\centering
\begin{tikzpicture}[
  font=\footnotesize,
  src/.style ={draw, rounded corners=2pt, align=center, text width=34mm,
               minimum height=13mm, inner sep=3pt},
  proc/.style={draw, rounded corners=2pt, align=center, text width=48mm,
               minimum height=16mm, inner sep=3pt},
  wide/.style={draw, rounded corners=2pt, align=center, text width=112mm,
               minimum height=9mm, inner sep=3pt},
  arr/.style ={-{Stealth[length=2mm]}, semithick}
]
  \node[src] (s1) at (-4.0, 2.3)
    {\textbf{Step 2}\\Privacy policies,\\store listings and labels};
  \node[src] (s2) at ( 0.0, 2.3)
    {\textbf{Step 2}\\Provider websites,\\ads, privacy blogs};
  \node[src] (s3) at ( 4.0, 2.3)
    {\textbf{Step 3}\\Manual privacy\\walkthrough (in-app)};

  \node[wide] (corpus) at (0,0)
    {\textbf{Frozen corpus} --- segmented units, each with source, language,
     app version and capture date};

  \node[proc] (llm) at (-3.0,-2.2)
    {\textbf{Step 4a}\\LLM-assisted coding,\\versioned prompts,
     JSON $+$ verbatim span};
  \node[proc] (man) at ( 3.0,-2.2)
    {\textbf{Step 4b}\\Manual coding of a\\stratified 20\,\% subsample};

  \node[wide] (adj) at (0,-4.3)
    {\textbf{Adjudicated codings} $+$ per-category agreement ($\kappa$)};
  \node[wide] (ana) at (0,-5.8)
    {\textbf{Analysis:} harm coverage (RQ1) $\cdot$ stance (RQ2) $\cdot$
     PET framing (RQ3) $\cdot$ cross-surface consistency (RQ4)};

  \draw[arr] (s1) -- (corpus);
  \draw[arr] (s2) -- (corpus);
  \draw[arr] (s3) -- (corpus);
  \draw[arr] (corpus) -- (llm);
  \draw[arr] (corpus) -- (man);
  \draw[arr] (llm) -- (adj);
  \draw[arr] (man) -- (adj);
  \draw[arr] (adj) -- (ana);
\end{tikzpicture}
\caption{Overview of the study design. Step~1 (sampling) determines which apps
enter the three collection streams at the top.}
\label{fig:method}
\end{figure}

\subsection{Step 1: App Sampling}
\label{sec:sampling}

\paragraph{Platforms and market.}
Google Play and the Apple App Store, German storefront, German device locale.
Both platforms are covered because their structured disclosures differ (Data
Safety section versus Privacy Nutrition Labels) and RQ4 depends on that
difference.

\paragraph{Search.}
A fixed, pre-registered keyword list is used on both stores, in German and
English, because German listings routinely use English terms. The seed list
covers the two domains and their common synonyms: \emph{Fitness, Fitness
Tracker, Schrittz\"ahler, Aktivit\"atstracker, Schlaf-Tracker, Herzfrequenz,
Kalorienz\"ahler, Ern\"ahrungstagebuch, Workout} for the fitness domain, and
\emph{Periode, Zyklus, Zyklustracker, Menstruation, Eisprung, Fruchtbarkeit,
Kinderwunsch, period tracker, ovulation, fertility} for the fertility domain.
The list is frozen before collection; any addition is recorded in the change log
of the sampling protocol with its date and reason. Store categories considered
are \emph{Health \& Fitness} and \emph{Medical}.

\paragraph{Operationalising \enquote{highest rated}.}
\enquote{Highest rated} is ambiguous and is therefore fixed in advance. For each
keyword, the first 20 results in the store's default ordering are recorded
together with their average rating and rating count. The pooled, de-duplicated
candidate set is then filtered to apps with at least 1{,}000 ratings on the
platform in question --- below that, the average rating is noise --- and ranked
by average rating, with rating count as the tie-break. This makes the sample
\emph{highly rated among widely used apps} rather than \emph{rated highly by a
handful of people}, which is the intended reading.

\paragraph{Inclusion criteria.} An app is included if all of the following hold:
\begin{enumerate}[leftmargin=*,itemsep=1pt]
  \item it is available in the German storefront of at least one of the two
        platforms;
  \item it collects health, fitness, or body-related self-tracking data from the
        user or a paired device;
  \item it is free to install, so that the walkthrough can reach onboarding
        without a purchase;
  \item a privacy policy is reachable from the store listing;
  \item its interface is available in German or English.
\end{enumerate}

\paragraph{Exclusion criteria.} An app is excluded if any of the following hold:
\begin{enumerate}[leftmargin=*,itemsep=1pt]
  \item no privacy policy is reachable, or the link is dead at capture time
        (recorded as a finding, then excluded from text coding);
  \item it is a prescription DiGA or is marketed as a regulated medical device;
  \item it is a pure device companion with no tracking function of its own;
  \item account creation requires identity verification or payment, which the
        ethics rules in Section~\ref{sec:ethics} forbid;
  \item it duplicates an app already in the sample from the same provider under
        a different name.
\end{enumerate}

\paragraph{Sample size and stratification.}
The target is \textbf{50 apps} for document analysis, stratified as 30 fitness
and activity apps and 20 fertility and period apps. The fertility stratum is
deliberately over-represented relative to its market share: it is the domain
where the harm literature is most developed and where the consequences are most
concrete~\parencite{cao2024deleted,malki2024exploring,shipp2020period}, which
makes it the more informative half of a comparison. A \textbf{pilot of five
apps} (three fitness, two fertility) is run first to test the collection
protocol, the codebook, and the prompts; pilot apps are re-coded with the final
instruments and retained in the sample.

\paragraph{Freezing.}
Everything is recorded with a capture date and app version, and the corpus is
frozen on a stated date. Policies change; a moving corpus cannot be analysed.
The full record per app is specified in
\pathname{04_Study/App-Sample/sampling_protocol.md}.

\subsection{Step 2: Material Collection}

For every app in the sample, the following are collected, archived as both the
original file and extracted plain text, and given a stable document identifier:

\begin{enumerate}[leftmargin=*,itemsep=2pt]
  \item \textbf{Privacy policy}, in German if one is offered to a German user and
        in English otherwise; both if both exist, with the difference recorded.
  \item \textbf{Store listing}: full description text, plus the structured
        disclosure (Play Data Safety section, App Store Privacy Nutrition Label)
        captured as text and screenshot.
  \item \textbf{Provider website}: the service description, any dedicated privacy
        or security page, and any privacy-related blog or newsroom post.
  \item \textbf{Advertisements}, where encountered --- in-app ads shown during
        the walkthrough, and any ad creative reachable from the provider's own
        channels. Advertising is treated as a \emph{secondary} corpus: it is
        coded when found, but its absence for a given app is not itself a
        finding, because ad delivery cannot be systematically reproduced.
  \item \textbf{In-app privacy text}, from the walkthrough (Step~3).
\end{enumerate}

Each document is recorded with source type, language, URL, capture date,
SHA-256 hash of the archived file, and the app version it belongs to. Collection
is manual or uses officially permitted access; each store's terms are checked
before any automated retrieval, and manual capture is preferred where the terms
are restrictive.

\paragraph{Segmentation.}
Documents are split into coding units at heading and paragraph boundaries, with
list items as their own units. Units, not documents, are the unit of coding;
document-level counts are derived afterwards. This keeps a 40-page policy and a
three-line store blurb commensurable.

\subsection{Step 3: Privacy-Focused Cognitive Walkthrough}
\label{sec:walkthrough}

Walkthroughs are manual and are run on a \textbf{stratified subsample of 12
apps}, drawn from the top of the ranked sample so that the most-used apps are
covered. Unlike the document sample, the subsample is \emph{balanced}: six
fitness and six fertility apps, each domain split three on Android and three on
iOS. Balance is deliberate here, because RQ4 compares domains and platforms
against each other and an unequal subsample would confound that comparison with
sample size. Twelve is what a single analyst can walk through consistently
within the schedule while covering both platforms and both domains.

The method follows the cognitive walkthrough~\parencite{wharton1994cognitive},
with one adaptation: the tasks are \emph{privacy} tasks. For each task the
analyst walks the action sequence and asks Wharton et al.'s four questions ---
will the user try to achieve the right effect; will they notice that the correct
action is available; will they associate the correct action with the effect they
want; and, once the action is performed, will they see that progress is being
made. The tasks are:

\begin{enumerate}[leftmargin=*,itemsep=1pt]
  \item find out what data the app collects;
  \item find out what the app or provider \emph{derives} from that data;
  \item find out what happens if the data is exposed, and how incidents are
        handled;
  \item reach the privacy policy from inside the app;
  \item turn off sharing with third parties;
  \item export the data;
  \item delete the account and the data.
\end{enumerate}

Tasks 2 and 3 are the ones that matter for RQ1: they are expected to fail, and
the walkthrough records \emph{where} they fail. Each session follows the fixed
sequence in \pathname{04_Study/Instrument/walkthrough_protocol.md} --- store
listing captured before install, fresh install, no restored backup, screen
recording throughout, every consent and permission dialogue recorded verbatim,
taps-to-target counted from the home screen. Observation and interpretation are
written separately in the notes; only the observation is data.

Walkthroughs are run by a single analyst, which is a known weakness of
inspection methods~\parencite{nielsen1994heuristic}. It is mitigated by the
fixed protocol, the screen recordings (which make a second pass possible), and
by double-coding the walkthrough notes as part of the validation subsample.

\subsection{Step 4: Coding Framework}
\label{sec:framework}

Every coding unit is coded on three axes. The axes are independent: a unit can
carry codes on all three, or on none.

\paragraph{Axis A --- harm type and pathway (RQ1).}
The backbone is the Citron and Solove typology: physical, economic,
reputational, psychological, autonomy, discrimination, and relationship
harm~\parencite{citron2021harms}. Its seven categories were built for
adjudication, not for content analysis, so each is given an operational
definition, a positive example, and a counter-example in the codebook; a
category that cannot be given a counter-example is not codable and is merged or
dropped. On top of this, four \textbf{probe codes} capture the specific silences
identified in Section~\ref{sec:problem}: \textsc{inference},
\textsc{breach-consequence}, \textsc{incident-handling}, and
\textsc{concrete-consequence}. Finally, where a unit does describe a route to
harm, it is annotated with the stage of the Knabenhans et al.\ harm
pathway~\parencite{knabenhans2026harms} it refers to --- system component,
adversarial capability (distinguishing, tracking, inferring), or harm --- which
records whether providers describe mechanisms, capabilities, or outcomes.

\paragraph{Axis B --- stance (RQ2).}
For each harm code, one stance is assigned: \textsc{acknowledged} (the harm is
named as a real possibility), \textsc{hedged} (named but qualified into
improbability), \textsc{minimised} (named alongside a claim that defuses it
without substantiation), \textsc{deflected} (named but made the user's
responsibility), or \textsc{absent} (recorded at document level, not unit
level). \textsc{Absent} is a document-level and app-level judgement: it can only
be asserted after the whole corpus for that app has been coded.

\paragraph{Axis C --- PET framing (RQ3).}
Where a protective technology or practice is described, the unit is coded for
which of the three principles of Martin et al.\ it invokes --- limiting
third-party access, minimising raw data in favour of inferences, or on-device
processing~\parencite{martin2025cookies} --- and for whether the protection is
framed as \textbf{secrecy} (nobody else sees the data) or as \textbf{harm
prevention} (a specific bad outcome is prevented). A separate code marks
\textsc{reassurance-without-substance}: claims such as \enquote{your data is
safe with us} that carry no verifiable commitment.

The framework is deductive at the start and inductive at the edges: codes that
the corpus demands but the framework does not anticipate are added through
thematic analysis~\parencite{braun2006using}, each one triggering a new codebook
version with its rationale recorded.

\subsection{Step 4a: LLM-Assisted Coding}

The corpus is too large to code entirely by hand --- 50 apps across four or five
source types, with individual policies commonly running to several thousand
words~\parencite{mcdonald2008cost} --- and LLMs have been evaluated for exactly
this kind of policy analysis at scale~\parencite{rodriguez2024llm}. Existing tooling is not reused directly:
the OPP-115 scheme and Polisis classify \emph{data
practices}~\parencite{wilson2016creation,harkous2018polisis}, and have no harm
categories, so a purpose-built prompt pipeline over the codebook of
Section~\ref{sec:framework} is required. Their category sets are, however, used
as a cross-check that the corpus was segmented and parsed correctly.

The pipeline is constrained as follows:

\begin{itemize}[leftmargin=*,itemsep=2pt]
  \item \textbf{One prompt per axis}, not one prompt for everything. Each prompt
        file carries the header specified in
        \pathname{04_Study/Prompts/README.md}: task, version, date, exact model
        name and version, input, output schema, and what changed since the
        previous version.
  \item \textbf{Structured output.} The model returns JSON: unit identifier,
        codes assigned, and for each code a \emph{verbatim span} copied from the
        unit plus a confidence flag. A code without a span is discarded
        automatically.
  \item \textbf{Determinism and stability.} A single pinned model version,
        temperature 0, three independent runs per unit, majority vote across
        runs. The rate of disagreement between runs is itself reported as a
        measure of instrument stability.
  \item \textbf{No inference beyond the text.} The prompt instructs the model to
        code what is written, never what the provider presumably means --- the
        same rule that binds the human coder.
  \item \textbf{Logging.} Every run is entered in
        \pathname{04_Study/Prompts/run_log.md} with date, prompt version, model,
        document count, and anything that went wrong.
\end{itemize}

\subsection{Step 4b: Validation}
\label{sec:validation}

LLM output is treated as a measurement instrument that must be calibrated, not
as ground truth.

\begin{enumerate}[leftmargin=*,itemsep=3pt]
  \item \textbf{Agreement.} A stratified 20\,\% subsample of coding units ---
        drawn across both domains, both platforms, and all source types --- is
        coded manually by the author, blind to the model output. Agreement is
        computed per code with Cohen's $\kappa$~\parencite{cohen1960coefficient}
        and reported per category, not as a single pooled figure, because the
        rare categories are the interesting ones and a pooled figure hides them.
  \item \textbf{Threshold and remedy.} The target is $\kappa \geq 0.70$ per
        code. Where a code falls below it, the failure is diagnosed as either a
        codebook problem (ambiguous definition, missing counter-example) or a
        prompt problem; the relevant instrument is revised, its version
        incremented, and a fresh subsample drawn. Codes that cannot be brought
        above the threshold after two revisions are reported as unreliable and
        excluded from quantitative claims, with the reason stated.
  \item \textbf{Precision audit.} Independently of agreement, every code that
        enters the results is checked against its verbatim span; a span that
        does not support the code is discarded and counted. The resulting false
        positive rate is reported.
  \item \textbf{Second coder.} A second coder on a portion of the subsample
        would strengthen the design considerably; whether one can be arranged
        through the supervisor is one of the open points in
        Section~\ref{sec:openpoints}. If none is available, this is stated
        plainly as a limitation and reporting follows the norms for
        single-coder qualitative work in HCI, which require the reliability
        procedure to be described rather than a number to be
        claimed~\parencite{mcdonald2019reliability}.
\end{enumerate}

\subsection{Codebook and Coding Procedure}

The codebook lives in \pathname{04_Study/Codebook/} and is versioned. Every
category carries a definition, a positive example, and an explicit
counter-example; the counter-example is mandatory because without it neither a
human coder nor a model has a boundary, and both drift. The standing rules are:

\begin{enumerate}[leftmargin=*,itemsep=1pt]
  \item Code the text as written, never what the provider is assumed to mean.
  \item A unit may carry several codes; codes are not mutually exclusive across
        axes.
  \item Every code records source, verbatim quote, and location.
  \item When unsure, code it and flag it --- never silently drop it.
  \item Any change to a category produces a new codebook version with date,
        change, and reason in the version history.
\end{enumerate}

Disagreements between the manual and model codings are adjudicated by the author
against the verbatim span and the codebook; each adjudication is logged with the
rule that resolved it, and a pattern of adjudications against the same category
is treated as evidence that the category needs revision.

\subsection{Ethics and Data Handling}
\label{sec:ethics}

The study involves \textbf{no human subjects}. The material is public documents
and the researcher's own interaction with publicly available apps, so no ethics
approval is expected to be required. This will nevertheless be confirmed in
writing with the supervisor and the relevant RUB\,/\,MPI-SP contact before
collection starts, and the answer filed in \pathname{04_Study/Ethics/}. If the
scenario-based survey is later added as a second phase, it does require review,
and an application would have to be filed well in advance.

The practical rules for the walkthroughs are binding regardless of that answer:
throwaway accounts and synthetic data only; no real name, address, health data,
or payment details entered into any studied app; no account that requires
identity verification (such apps are excluded by
Section~\ref{sec:sampling}); each store's terms checked before any automated
collection. Screen recordings and archived documents are stored locally in
\pathname{05_Data/raw/}, which is excluded from version control; only derived,
non-identifying data is committed. Use of an LLM in the analysis is part of the
method and is declared as such, both in the methodology chapter and in the
statutory declaration on AI use.

\subsection{Limitations and Threats to Validity}

\paragraph{Absence is not intent.}
The study can establish that a harm type is not mentioned. It cannot establish
why. Silence may be a drafting convention, a legal risk decision, or a
deliberate omission, and the thesis will not attribute motive. This is a
boundary on the claims, not a flaw in the measurement.

\paragraph{A snapshot, not a trend.}
Policies, labels, and interfaces change continuously. The corpus is frozen on a
stated date and every claim is indexed to it. Longitudinal comparison is out of
scope.

\paragraph{Market and language.}
The German storefront determines the sample. Many providers operate
internationally and serve near-identical policies everywhere, but jurisdiction
shapes the text, and findings should not be read as global without
qualification. Where a policy exists in both German and English, the pair is
recorded, and any divergence between them is itself reportable.

\paragraph{Measurement error in the instrument.}
LLM coding is sensitive to prompt wording and is not perfectly reproducible even
at temperature 0. The mitigations are in Section~\ref{sec:validation}: pinned
model version, repeated runs with a reported disagreement rate, per-category
$\kappa$ against manual coding, mandatory verbatim spans, and a precision audit.
The agreement figures are reported whatever they turn out to be.

\paragraph{Single analyst.}
One person builds the sample, runs the walkthroughs, and does the manual coding.
Fixed protocols, screen recordings, versioned instruments, and a committed audit
trail are the compensating controls; a second coder, if one can be arranged,
would be better and is requested in Section~\ref{sec:openpoints}.

\paragraph{Coverage of the interface.}
Walkthroughs cover 12 of 50 apps and only the free tier. Privacy-relevant
content behind a paywall or a rarely used flow may be missed. The subsample is
drawn from the top of the ranking so that the apps most people actually use are
the ones inspected.

\paragraph{Advertising is opportunistic.}
Ad delivery is personalised and cannot be reproduced systematically, so the
advertising corpus is a secondary source. Presence of privacy claims in ads is
reportable; their absence is not.

\paragraph{The text is not the practice.}
Coding communication says nothing about whether the communication is true.
That question is answered by traffic and code analysis in existing
work~\parencite{arkalakis2024abandon,hudig2025intimate,huckvale2019assessment},
and the findings here should be read alongside it, not as a substitute for it.
